\section{Conclusões}
\label{sec:conclusoes}

Este trabalho desenvolveu um sistema integrado de aquisição para a bancada de
ensaios de conjuntos motopropulsores da equipe Albatroz Aerodesign, reunindo
em uma placa de medição a aquisição das grandezas elétricas do motor, a
estimação de rotação e torque por um observador por modos deslizantes, o
acionamento do motor e a segurança, e em uma aplicação de computador a
supervisão e o registro dos ensaios. Os objetivos específicos são retomados a
seguir.

O sistema de medição do motor foi projetado e construído: a placa de
condicionamento das fases leva as três tensões e as correntes de duas fases,
medidas por sensores de efeito Hall com núcleo toroidal, aos dois conversores
do microcontrolador, amostradas aos pares no mesmo instante, a
\SI{111}{\kilo\hertz}. O conjunto núcleo--sensor foi caracterizado em bancada,
e as cadeias de tensão e corrente foram dimensionadas para a faixa do
conversor, com filtros de mesma frequência de corte nos dois ramos. Nos
ensaios com o motor em rotação, as tensões lidas pela placa alcançaram a do
barramento nas três fases e as correntes em repouso ficaram próximas de zero
após o levantamento automático do deslocamento.

O motor foi modelado e os parâmetros necessários ao observador foram
identificados: resistência de fase de \SI{8,69}{\milli\ohm}, indutância pelo
método da ressonância, perdas a vazio $B$ e $\tau_c$ por regressão sobre
\num{19} patamares e constante de força contraeletromotriz equivalente a
\num{310}~RPM/V, próxima da especificação de placa. O observador por modos
deslizantes, com o PLL e a reconstrução analítica da força contraeletromotriz,
foi validado em simulação, com ganhos obtidos por ajuste automático, e aplicado
aos registros de bancada: a velocidade estimada teve erro médio de
\SI{0,092}{\percent}, o ângulo ficou em média dentro de \SI{0,3}{\degree} da
referência e o torque, calculado com a corrente medida, ficou
\SI{7,2}{\percent} abaixo do medido. Abaixo de cerca de \num{1930}~RPM o
sincronismo não trava, o que limita a estimação de torque nessa faixa.

A medição, o acionamento e a segurança foram integrados na placa de medição.
O observador roda embarcado, um passo a cada amostra, e o processamento de
cada lote de amostras termina antes da chegada do seguinte, sem descarte nos
ensaios realizados. A placa comanda o controlador do motor e os relés de
potência com pré-carga, e a parada de emergência atua por meio físico,
independentemente do microcontrolador. Com as hélices APC 15x6E e 17x10E, a
rotação e o torque estimados assentaram em patamares a cada degrau de
acelerador e reproduziram a forma das curvas do fabricante, com valores acima
delas, diferença discutida na \autoref{sec:discussao_helices}.

A placa e a aplicação de supervisão foram interligadas por rede Ethernet, com
troca de pacotes por UDP. Comandos e telemetria são descritos em um único
arquivo de configuração, do qual são gerados os cabeçalhos do
\textit{firmware} e que a aplicação lê para montar e interpretar os pacotes,
o que mantém as duas pontas coerentes. A telemetria chega ao computador a
\SI{10}{\hertz}, com as grandezas medidas, as estimadas e os tempos de
processamento da placa.

Por fim, a aplicação de supervisão permite monitorar e comandar a bancada,
calibrar os sensores, ajustar os parâmetros do observador, executar os
procedimentos de estimação e registrar os ensaios em CSV. Os ajustes feitos
pela interface ficam gravados na memória não volátil da placa.

Em conjunto, os resultados indicam que o sistema cumpre o objetivo geral de
adaptar a bancada para ensaios instrumentados de conjuntos motopropulsores,
permitindo comparar as grandezas medidas com as estimadas pelo observador. A
principal limitação é a ausência de uma medida de torque independente no
mesmo eixo, que permitiria quantificar diretamente o erro da estimativa de
torque nos ensaios das hélices.
